\IfFileExists{stacks-project.cls}{%
\documentclass{stacks-project}
}{%
\documentclass{amsart}
}

% For dealing with references we use the comment environment
\usepackage{verbatim}
\newenvironment{reference}{\comment}{\endcomment}
%\newenvironment{reference}{}{}
\newenvironment{slogan}{\comment}{\endcomment}
\newenvironment{history}{\comment}{\endcomment}

% For commutative diagrams we use Xy-pic
\usepackage[all]{xy}

% We use 2cell for 2-commutative diagrams.
\xyoption{2cell}
\UseAllTwocells

% We use multicol for the list of chapters between chapters
\usepackage{multicol}

% This is generally recommended for better output
\usepackage{lmodern}
\usepackage[T1]{fontenc}

% For cross-file-references
\usepackage{xr-hyper}

% Package for hypertext links:
\usepackage{hyperref}

% For any local file, say "hello.tex" you want to link to please
% use \externaldocument[hello-]{hello}
\externaldocument[introduction-]{introduction}
\externaldocument[conventions-]{conventions}
\externaldocument[sets-]{sets}
\externaldocument[categories-]{categories}
\externaldocument[topology-]{topology}
\externaldocument[sheaves-]{sheaves}
\externaldocument[sites-]{sites}
\externaldocument[stacks-]{stacks}
\externaldocument[fields-]{fields}
\externaldocument[algebra-]{algebra}
\externaldocument[brauer-]{brauer}
\externaldocument[homology-]{homology}
\externaldocument[derived-]{derived}
\externaldocument[simplicial-]{simplicial}
\externaldocument[more-algebra-]{more-algebra}
\externaldocument[smoothing-]{smoothing}
\externaldocument[modules-]{modules}
\externaldocument[sites-modules-]{sites-modules}
\externaldocument[injectives-]{injectives}
\externaldocument[cohomology-]{cohomology}
\externaldocument[sites-cohomology-]{sites-cohomology}
\externaldocument[dga-]{dga}
\externaldocument[dpa-]{dpa}
\externaldocument[sdga-]{sdga}
\externaldocument[hypercovering-]{hypercovering}
\externaldocument[schemes-]{schemes}
\externaldocument[constructions-]{constructions}
\externaldocument[properties-]{properties}
\externaldocument[morphisms-]{morphisms}
\externaldocument[coherent-]{coherent}
\externaldocument[divisors-]{divisors}
\externaldocument[limits-]{limits}
\externaldocument[varieties-]{varieties}
\externaldocument[topologies-]{topologies}
\externaldocument[descent-]{descent}
\externaldocument[perfect-]{perfect}
\externaldocument[more-morphisms-]{more-morphisms}
\externaldocument[flat-]{flat}
\externaldocument[groupoids-]{groupoids}
\externaldocument[more-groupoids-]{more-groupoids}
\externaldocument[etale-]{etale}
\externaldocument[chow-]{chow}
\externaldocument[intersection-]{intersection}
\externaldocument[pic-]{pic}
\externaldocument[weil-]{weil}
\externaldocument[adequate-]{adequate}
\externaldocument[dualizing-]{dualizing}
\externaldocument[duality-]{duality}
\externaldocument[discriminant-]{discriminant}
\externaldocument[derham-]{derham}
\externaldocument[local-cohomology-]{local-cohomology}
\externaldocument[algebraization-]{algebraization}
\externaldocument[curves-]{curves}
\externaldocument[resolve-]{resolve}
\externaldocument[models-]{models}
\externaldocument[functors-]{functors}
\externaldocument[equiv-]{equiv}
\externaldocument[pione-]{pione}
\externaldocument[etale-cohomology-]{etale-cohomology}
\externaldocument[proetale-]{proetale}
\externaldocument[relative-cycles-]{relative-cycles}
\externaldocument[more-etale-]{more-etale}
\externaldocument[trace-]{trace}
\externaldocument[crystalline-]{crystalline}
\externaldocument[spaces-]{spaces}
\externaldocument[spaces-properties-]{spaces-properties}
\externaldocument[spaces-morphisms-]{spaces-morphisms}
\externaldocument[decent-spaces-]{decent-spaces}
\externaldocument[spaces-cohomology-]{spaces-cohomology}
\externaldocument[spaces-limits-]{spaces-limits}
\externaldocument[spaces-divisors-]{spaces-divisors}
\externaldocument[spaces-over-fields-]{spaces-over-fields}
\externaldocument[spaces-topologies-]{spaces-topologies}
\externaldocument[spaces-descent-]{spaces-descent}
\externaldocument[spaces-perfect-]{spaces-perfect}
\externaldocument[spaces-more-morphisms-]{spaces-more-morphisms}
\externaldocument[spaces-flat-]{spaces-flat}
\externaldocument[spaces-groupoids-]{spaces-groupoids}
\externaldocument[spaces-more-groupoids-]{spaces-more-groupoids}
\externaldocument[bootstrap-]{bootstrap}
\externaldocument[spaces-pushouts-]{spaces-pushouts}
\externaldocument[spaces-chow-]{spaces-chow}
\externaldocument[groupoids-quotients-]{groupoids-quotients}
\externaldocument[spaces-more-cohomology-]{spaces-more-cohomology}
\externaldocument[spaces-simplicial-]{spaces-simplicial}
\externaldocument[spaces-duality-]{spaces-duality}
\externaldocument[formal-spaces-]{formal-spaces}
\externaldocument[restricted-]{restricted}
\externaldocument[spaces-resolve-]{spaces-resolve}
\externaldocument[formal-defos-]{formal-defos}
\externaldocument[defos-]{defos}
\externaldocument[cotangent-]{cotangent}
\externaldocument[examples-defos-]{examples-defos}
\externaldocument[algebraic-]{algebraic}
\externaldocument[examples-stacks-]{examples-stacks}
\externaldocument[stacks-sheaves-]{stacks-sheaves}
\externaldocument[criteria-]{criteria}
\externaldocument[artin-]{artin}
\externaldocument[quot-]{quot}
\externaldocument[stacks-properties-]{stacks-properties}
\externaldocument[stacks-morphisms-]{stacks-morphisms}
\externaldocument[stacks-limits-]{stacks-limits}
\externaldocument[stacks-cohomology-]{stacks-cohomology}
\externaldocument[stacks-perfect-]{stacks-perfect}
\externaldocument[stacks-introduction-]{stacks-introduction}
\externaldocument[stacks-more-morphisms-]{stacks-more-morphisms}
\externaldocument[stacks-geometry-]{stacks-geometry}
\externaldocument[moduli-]{moduli}
\externaldocument[moduli-curves-]{moduli-curves}
\externaldocument[examples-]{examples}
\externaldocument[exercises-]{exercises}
\externaldocument[guide-]{guide}
\externaldocument[desirables-]{desirables}
\externaldocument[coding-]{coding}
\externaldocument[obsolete-]{obsolete}
\externaldocument[fdl-]{fdl}
\externaldocument[index-]{index}

% Theorem environments.
%
\theoremstyle{plain}
\newtheorem{theorem}[subsection]{Theorem}
\newtheorem{proposition}[subsection]{Proposition}
\newtheorem{lemma}[subsection]{Lemma}

\theoremstyle{definition}
\newtheorem{definition}[subsection]{Definition}
\newtheorem{example}[subsection]{Example}
\newtheorem{exercise}[subsection]{Exercise}
\newtheorem{situation}[subsection]{Situation}

\theoremstyle{remark}
\newtheorem{remark}[subsection]{Remark}
\newtheorem{remarks}[subsection]{Remarks}

\numberwithin{equation}{subsection}

% Macros
%
\def\lim{\mathop{\mathrm{lim}}\nolimits}
\def\colim{\mathop{\mathrm{colim}}\nolimits}
\def\Spec{\mathop{\mathrm{Spec}}}
\def\Hom{\mathop{\mathrm{Hom}}\nolimits}
\def\Ext{\mathop{\mathrm{Ext}}\nolimits}
\def\SheafHom{\mathop{\mathcal{H}\!\mathit{om}}\nolimits}
\def\SheafExt{\mathop{\mathcal{E}\!\mathit{xt}}\nolimits}
\def\Sch{\mathit{Sch}}
\def\Mor{\mathop{\mathrm{Mor}}\nolimits}
\def\Ob{\mathop{\mathrm{Ob}}\nolimits}
\def\Sh{\mathop{\mathit{Sh}}\nolimits}
\def\NL{\mathop{N\!L}\nolimits}
\def\CH{\mathop{\mathrm{CH}}\nolimits}
\def\proetale{{pro\text{-}\acute{e}tale}}
\def\etale{{\acute{e}tale}}
\def\QCoh{\mathit{QCoh}}
\def\Ker{\mathop{\mathrm{Ker}}}
\def\Im{\mathop{\mathrm{Im}}}
\def\Coker{\mathop{\mathrm{Coker}}}
\def\Coim{\mathop{\mathrm{Coim}}}

% Boxtimes
%
\DeclareMathSymbol{\boxtimes}{\mathbin}{AMSa}{"02}

%
% Macros for moduli stacks/spaces
%
\def\QCohstack{\mathcal{QC}\!\mathit{oh}}
\def\Cohstack{\mathcal{C}\!\mathit{oh}}
\def\Spacesstack{\mathcal{S}\!\mathit{paces}}
\def\Quotfunctor{\mathrm{Quot}}
\def\Hilbfunctor{\mathrm{Hilb}}
\def\Curvesstack{\mathcal{C}\!\mathit{urves}}
\def\Polarizedstack{\mathcal{P}\!\mathit{olarized}}
\def\Complexesstack{\mathcal{C}\!\mathit{omplexes}}
% \Pic is the operator that assigns to X its picard group, usage \Pic(X)
% \Picardstack_{X/B} denotes the Picard stack of X over B
% \Picardfunctor_{X/B} denotes the Picard functor of X over B
\def\Pic{\mathop{\mathrm{Pic}}\nolimits}
\def\Picardstack{\mathcal{P}\!\mathit{ic}}
\def\Picardfunctor{\mathrm{Pic}}
\def\Deformationcategory{\mathcal{D}\!\mathit{ef}}
\usepackage{cleveref}
\crefname{originalreference}{original reference}{original references}
\Crefname{originalreference}{Original reference}{Original references}
\makeatletter
\newlabel{groupoids-lemma-quotient-groupoid-restrict}{{02VH}{source}{}{sourcepointer-1}{}}
\newlabel{groupoids-lemma-quotient-groupoid-restrict@cref}{{[originalreference][0][]02VH}{[1][0][]source}{}{}{}}
\makeatother
\crefname{originalremark}{original remark}{original remarks}
\Crefname{originalremark}{Original remark}{Original remarks}
\makeatletter
\newlabel{remark-recall}{{03II}{source}{}{sourcepointer-2}{}}
\newlabel{remark-recall@cref}{{[originalremark][0][]03II}{[1][0][]source}{}{}{}}
\makeatother
\crefname{originalreference}{original reference}{original references}
\Crefname{originalreference}{Original reference}{Original references}
\makeatletter
\newlabel{schemes-lemma-injective-points-surjective-stalks}{{01L6}{source}{}{sourcepointer-3}{}}
\newlabel{schemes-lemma-injective-points-surjective-stalks@cref}{{[originalreference][0][]01L6}{[1][0][]source}{}{}{}}
\makeatother
\crefname{originalreference}{original reference}{original references}
\Crefname{originalreference}{Original reference}{Original references}
\makeatletter
\newlabel{schemes-lemma-mono-towards-spec-field}{{03DP}{source}{}{sourcepointer-4}{}}
\newlabel{schemes-lemma-mono-towards-spec-field@cref}{{[originalreference][0][]03DP}{[1][0][]source}{}{}{}}
\makeatother
\crefname{originalreference}{original reference}{original references}
\Crefname{originalreference}{Original reference}{Original references}
\makeatletter
\newlabel{spaces-properties-lemma-point-like-spaces}{{03DZ}{source}{}{sourcepointer-5}{}}
\newlabel{spaces-properties-lemma-point-like-spaces@cref}{{[originalreference][0][]03DZ}{[1][0][]source}{}{}{}}
\makeatother
\crefname{originalreference}{original reference}{original references}
\Crefname{originalreference}{Original reference}{Original references}
\makeatletter
\newlabel{spaces-properties-lemma-points-cartesian}{{03H4}{source}{}{sourcepointer-6}{}}
\newlabel{spaces-properties-lemma-points-cartesian@cref}{{[originalreference][0][]03H4}{[1][0][]source}{}{}{}}
\makeatother
\crefname{originalreference}{original reference}{original references}
\Crefname{originalreference}{Original reference}{Original references}
\makeatletter
\newlabel{spaces-theorem-presentation}{{02WW}{source}{}{sourcepointer-7}{}}
\newlabel{spaces-theorem-presentation@cref}{{[originalreference][0][]02WW}{[1][0][]source}{}{}{}}
\makeatother
\title{Four equivalent characterizations of finiteness of the groupoid fibre over a point}
\author{The Stacks Project Authors}
\date{stacks-project:decent-spaces.tex:local-snapshot-20261012; GNU FDL 1.2 or later}
\begin{document}
\maketitle
\noindent\textit{Editorial note.} No omission markers in-span; every direction ends with an explicit conclusion.\par\medskip
\noindent\textit{Allowed standard prerequisites (editorial).}\begin{itemize}
\item Monomorphisms from points. Schemes, lemma injective-points-surjective-stalks and lemma mono-towards-spec-field (cross-chapter, named). Identify points with monomorphisms and recognize finite disjoint unions of spectra of fields.
\item Points lift through etale morphisms. Properties of Spaces, lemma points-cartesian (cross-chapter, named). Provides the lifting points w-tilde, w'-tilde and the final fibre identifications.
\item Presentations give algebraic spaces. Spaces, theorem presentation (cross-chapter, named). The quotient X' = Spec(k)/R is an algebraic space in (3) implies (4).
\item Point-like spaces. Properties of Spaces, lemma point-like-spaces (cross-chapter, named). Identifies X' with Spec(k') in (3) implies (4).
\end{itemize}\medskip
\ifdefined\hypertarget\hypertarget{sourcepointer-1}{}\fi
\noindent\textit{Source reference.} Original reference 02VH: focused-seeds/geometry/S9-decent-r-finite-above-x/ref-tags.txt, source lines 5--5. Navigation pointer only: the counted proof cites this statement by \textbackslash{}ref as a standing input; the deduction the counted proof has to make is carried out inside the counted span, and no part of that deduction is delegated to this pointer.\par\medskip
\ifdefined\hypertarget\hypertarget{sourcepointer-2}{}\fi
\noindent\textit{Source reference.} Original remark 03II: sources/stacks-decent-spaces/pinned-original-source.tex, source lines 459--459. Navigation pointer only: the counted proof cites this statement by \textbackslash{}ref as a standing input; the deduction the counted proof has to make is carried out inside the counted span, and no part of that deduction is delegated to this pointer.\par\medskip
\ifdefined\hypertarget\hypertarget{sourcepointer-3}{}\fi
\noindent\textit{Source reference.} Original reference 01L6: focused-seeds/geometry/S9-decent-r-finite-above-x/ref-tags.txt, source lines 6--6. Navigation pointer only: the counted proof cites this statement by \textbackslash{}ref as a standing input; the deduction the counted proof has to make is carried out inside the counted span, and no part of that deduction is delegated to this pointer.\par\medskip
\ifdefined\hypertarget\hypertarget{sourcepointer-4}{}\fi
\noindent\textit{Source reference.} Original reference 03DP: focused-seeds/geometry/S9-decent-r-finite-above-x/ref-tags.txt, source lines 7--7. Navigation pointer only: the counted proof cites this statement by \textbackslash{}ref as a standing input; the deduction the counted proof has to make is carried out inside the counted span, and no part of that deduction is delegated to this pointer.\par\medskip
\ifdefined\hypertarget\hypertarget{sourcepointer-5}{}\fi
\noindent\textit{Source reference.} Original reference 03DZ: focused-seeds/geometry/S9-decent-r-finite-above-x/ref-tags.txt, source lines 8--8. Navigation pointer only: the counted proof cites this statement by \textbackslash{}ref as a standing input; the deduction the counted proof has to make is carried out inside the counted span, and no part of that deduction is delegated to this pointer.\par\medskip
\ifdefined\hypertarget\hypertarget{sourcepointer-6}{}\fi
\noindent\textit{Source reference.} Original reference 03H4: focused-seeds/geometry/S9-decent-r-finite-above-x/ref-tags.txt, source lines 9--9. Navigation pointer only: the counted proof cites this statement by \textbackslash{}ref as a standing input; the deduction the counted proof has to make is carried out inside the counted span, and no part of that deduction is delegated to this pointer.\par\medskip
\ifdefined\hypertarget\hypertarget{sourcepointer-7}{}\fi
\noindent\textit{Source reference.} Original reference 02WW: focused-seeds/geometry/S9-decent-r-finite-above-x/ref-tags.txt, source lines 10--10. Navigation pointer only: the counted proof cites this statement by \textbackslash{}ref as a standing input; the deduction the counted proof has to make is carried out inside the counted span, and no part of that deduction is delegated to this pointer.\par\medskip
\begin{lemma}
\label{lemma-R-finite-above-x}
Let $S$ be a scheme. Let $X$ be an algebraic space over $S$.
Let $x \in |X|$. The following are equivalent:
\begin{enumerate}
\item there exists a scheme $U$, an \'etale morphism
$\varphi : U \to X$, and points $u, u' \in U$ mapping to
$x$ such that setting $R = U \times_X U$ the fibre of
$$
|R| \to |U| \times_{|X|} |U|
$$
over $(u, u')$ is finite,
\item for every scheme $U$, \'etale morphism $\varphi : U \to X$ and
any points $u, u' \in U$ mapping to
$x$ setting $R = U \times_X U$ the fibre of
$$
|R| \to |U| \times_{|X|} |U|
$$
over $(u, u')$ is finite,
\item there exists a morphism $\Spec(k) \to X$ with $k$ a field
in the equivalence class of $x$ such that the projections
$\Spec(k) \times_X \Spec(k) \to \Spec(k)$ are
\'etale and quasi-compact, and
\item there exists a monomorphism $\Spec(k) \to X$ with $k$ a field
in the equivalence class of $x$.
\end{enumerate}
\end{lemma}

\begin{proof}
Assume (1), i.e., let $\varphi : U \to X$ be an \'etale morphism from a scheme
towards $X$, and let $u, u'$ be points of $U$ lying over $x$
such that the fibre of $|R| \to |U| \times_{|X|} |U|$ over $(u, u')$
is a finite set. In this proof we think of a point $u = \Spec(\kappa(u))$
as a scheme. Note that $u \to U$, $u' \to U$ are monomorphisms (see
Schemes, Lemma \ref{schemes-lemma-injective-points-surjective-stalks}),
hence $u \times_X u' \to R = U \times_X U$ is a monomorphism.
In this language the assumption really means that
$u \times_X u'$ is a scheme whose underlying topological space has
finitely many points.
Let $\psi : W \to X$ be an \'etale morphism from a scheme towards $X$.
Let $w, w' \in W$ be points of $W$ mapping to $x$.
We have to show that $w \times_X w'$ is a scheme whose underlying topological
space has finitely many points.
Consider the fibre product diagram
$$
\xymatrix{
W \times_X U \ar[r]_p \ar[d]_q & U \ar[d]^\varphi \\
W \ar[r]^\psi & X
}
$$
As $x$ is the image of $u$ and $u'$ we may pick points
$\tilde w, \tilde w'$ in $W \times_X U$ with $q(\tilde w) = w$,
$q(\tilde w') = w'$, $u = p(\tilde w)$ and $u' = p(\tilde w')$, see
Properties of Spaces,
Lemma \ref{spaces-properties-lemma-points-cartesian}.
As $p$, $q$ are \'etale the field extensions
$\kappa(w) \subset \kappa(\tilde w) \supset \kappa(u)$ and
$\kappa(w') \subset \kappa(\tilde w') \supset \kappa(u')$ are
finite separable, see Remark \ref{remark-recall}.
Then we get a commutative diagram
$$
\xymatrix{
w \times_X w' \ar[d] &
\tilde w \times_X \tilde w' \ar[l] \ar[d] \ar[r] &
u \times_X u' \ar[d] \\
w \times_X w' &
\tilde w \times_S \tilde w' \ar[l] \ar[r] &
u \times_S u'
}
$$
where the squares are fibre product squares. The lower horizontal
morphisms are \'etale and quasi-compact, as any scheme of the form
$\Spec(k) \times_S \Spec(k')$ is affine, and by our
observations about the field extensions above.
Thus we see that the top horizontal arrows are \'etale and quasi-compact
and hence have finite fibres.
We have seen above that $|u \times_X u'|$ is finite, so we conclude that
$|w \times_X w'|$ is finite. In other words, (2) holds.

\medskip\noindent
Assume (2). Let $U \to X$ be an \'etale morphism from a scheme $U$
such that $x$ is in the image of $|U| \to |X|$. Let $u \in U$ be
a point mapping to $x$. Then we have seen in the previous
paragraph that $u = \Spec(\kappa(u)) \to X$ has the property that
$u \times_X u$ has a finite underlying topological space. On the other
hand, the projection maps $u \times_X u \to u$ are the composition
$$
u \times_X u \longrightarrow
u \times_X U \longrightarrow
u \times_X X = u,
$$
i.e., the composition of a monomorphism (the base change of the monomorphism
$u \to U$) by an \'etale morphism (the base change of the \'etale morphism
$U \to X$). Hence $u \times_X U$ is a disjoint union of spectra of fields
finite separable over $\kappa(u)$ (see
Remark \ref{remark-recall}). Since $u \times_X u$ is finite the image
of it in $u \times_X U$ is a finite disjoint union of spectra of fields
finite separable over $\kappa(u)$. By
Schemes, Lemma \ref{schemes-lemma-mono-towards-spec-field}
we conclude that $u \times_X u$ is a finite disjoint union of spectra
of fields finite separable over $\kappa(u)$. In other words, we see that
$u \times_X u \to u$ is quasi-compact and \'etale. This means that (3) holds.

\medskip\noindent
Let us prove that (3) implies (4). Let $\Spec(k) \to X$ be a morphism
from the spectrum of a field into $X$, in the equivalence class of $x$
such that the two projections
$t, s : R = \Spec(k) \times_X \Spec(k)  \to \Spec(k)$
are quasi-compact and \'etale.
This means in particular
that $R$ is an \'etale equivalence relation on $\Spec(k)$.
By Spaces, Theorem \ref{spaces-theorem-presentation}
we know that the quotient sheaf
$X' = \Spec(k)/R$ is an algebraic space. By
Groupoids, Lemma \ref{groupoids-lemma-quotient-groupoid-restrict}
the map $X' \to X$ is a monomorphism.
Since $s, t$ are quasi-compact, we see that $R$ is quasi-compact and hence
Properties of Spaces,
Lemma \ref{spaces-properties-lemma-point-like-spaces}
applies to $X'$, and we see that
$X' = \Spec(k')$ for some field $k'$. Hence we get a factorization
$$
\Spec(k) \longrightarrow
\Spec(k') \longrightarrow X
$$
which shows that $\Spec(k') \to X$ is a monomorphism mapping
to $x \in |X|$. In other words (4) holds.

\medskip\noindent
Finally, we prove that (4) implies (1). Let $\Spec(k) \to X$
be a monomorphism with $k$ a field in the equivalence class of $x$.
Let $U \to X$ be a surjective \'etale morphism from a scheme $U$ to $X$.
Let $u \in U$ be a point over $x$. Since $\Spec(k) \times_X u$
is nonempty, and since $\Spec(k) \times_X u \to u$ is a monomorphism
we conclude that $\Spec(k) \times_X u = u$ (see
Schemes, Lemma \ref{schemes-lemma-mono-towards-spec-field}).
Hence $u \to U \to X$ factors through $\Spec(k) \to X$, here is
a picture
$$
\xymatrix{
u \ar[r] \ar[d] & U \ar[d] \\
\Spec(k) \ar[r] & X
}
$$
Since the right vertical arrow is \'etale this implies that
$\kappa(u)/k$ is a finite separable extension. Hence we conclude that
$$
u \times_X u = u \times_{\Spec(k)} u
$$
is a finite scheme, and we win by the discussion of the meaning of property
(1) in the first paragraph of this proof.
\end{proof}

\end{document}
